\documentclass[11pt]{article}
\usepackage[margin=1in]{geometry}
\usepackage{amsmath,amssymb}
\usepackage{array,booktabs}
\usepackage{tikz}
\usepackage{enumitem}

\setlength{\parindent}{0pt}
\setlength{\parskip}{6pt}
\begin{document}

\section*{Problem 8}

\textbf{Question.} If a random sample of size $15$ is taken from a Normal population with mean $75$, what are the chances that
\[
\left(\frac{\bar X-75}{S/\sqrt{15}}\right)
\]
exceeds $2.145$ in absolute value?

\textbf{Solution.} Since the population standard deviation is not given, use the sample standard deviation $S$. Then
\[
T=\frac{\bar X-\mu}{S/\sqrt{n}}
\]
has a Student's $t$-distribution with
\[
n-1=15-1=14
\]
degrees of freedom. Here,
\[
T=\frac{\bar X-75}{S/\sqrt{15}}\sim t_{14}.
\]
Find
\[
P\left(\left|\frac{\bar X-75}{S/\sqrt{15}}\right|>2.145\right)=P(|T|>2.145).
\]
From the $t$-table with $14$ degrees of freedom,
\[
t_{0.025,14}=2.145.
\]
Therefore,
\[
P(T>2.145)=0.025.
\]
By symmetry of the $t$-distribution,
\[
P(|T|>2.145)=2P(T>2.145)=2(0.025)=0.05.
\]
Hence,
\[
\boxed{P\left(\left|\frac{\bar X-75}{S/\sqrt{15}}\right|>2.145\right)=0.05}.
\]

\newpage

\end{document}
